\documentclass{article}

\usepackage[utf8]{inputenc}
\usepackage[T1]{fontenc}
\usepackage{helvet}
\usepackage{geometry}
\usepackage{xcolor}
\usepackage{eso-pic}
\usepackage{fancyhdr}
\usepackage{parskip}
\usepackage{microtype}
\usepackage[english, ukrainian]{babel}
\usepackage{url}
\usepackage{amsmath}
\usepackage{amssymb}
\usepackage{amsthm}
\usepackage{longtable}
\usepackage{booktabs}
\usepackage{hyperref}
\usepackage{titlesec}
\usepackage{graphicx}

\definecolor{nato}{HTML}{293757}
\definecolor{footergray}{gray}{0.65}

\geometry{a4paper,left=3cm,right=3cm,top=3cm,bottom=3cm}
\renewcommand{\familydefault}{\sfdefault}
\setcounter{secnumdepth}{3}

\titleformat{\section}
  {\color{nato}\Huge\bfseries}{\thesection.}{1em}{}
\titlespacing*{\section}{0pt}{30pt}{12pt}

\titleformat{\subsection}
  {\color{nato}\LARGE\bfseries}{\thesubsection.}{1em}{}
\titlespacing*{\subsection}{0pt}{20pt}{8pt}

\titleformat{\subsubsection}
  {\color{nato}\Large\bfseries}{\thesubsubsection.}{1em}{}
\titlespacing*{\subsubsection}{0pt}{10pt}{6pt}

\fancyhf{}
\fancyhead[R]{\small\color{footergray} 31 березня 2026}
\fancyfoot[L]{\small\color{footergray} Технічні вимоги ЄСІКС. V. Судопис}
\fancyfoot[R]{\small\color{footergray}\thepage}

\newcommand\HUGE[1]{\fontsize{40}{40}\selectfont #1}

\renewcommand{\headrulewidth}{0pt}
\renewcommand{\footrulewidth}{0.4pt}
\renewcommand{\footrule}{\color{footergray}\hrule width \textwidth height 0.4pt}

\begin{document}

\pagestyle{fancy}

\begin{titlepage}
  \AddToShipoutPictureBG*{\AtPageLowerLeft{\color{nato}\rule{\paperwidth}{\paperheight}}}
  \color{white}\thispagestyle{empty}\vspace*{4cm}\centering
  {\Huge  \textbf{Технічні вимоги ЄСІКС} \par} \vspace{2cm}
  {\HUGE  \textbf{V. Судопис} \par} \vspace{2.5cm}
  {\LARGE \textbf{V-1. Справи} \par}
  {\LARGE \textbf{V-2. Провадження} \par}
  {\LARGE \textbf{V-3. Судові засідання і їх рішення} \par}
  {\LARGE \textbf{V-4. Авторозподіл справ} \par} \vspace{2.5cm}
  \vfill
  {\large 2026 \copyright\ Інформаційні Судові Системи \par}
\end{titlepage}

\newpage

\begin{center}
  {\Huge\color{nato}\bfseries Зміст\par}
\end{center}
\vspace{40pt}
\tableofcontents

\begin{abstract}
ТЕХНІЧНІ ВИМОГИ
НА СТВОРЕННЯ ЗАСОБУ ІНФОРМАТИЗАЦІЇ - КОМПОНЕНТІВ “СУДОПИСУ” ЄДИНОЇ СУДОВОЇ ІНФОРМАЦІЙНО-КОМУНІКАЦІЙНОЇ СИСТЕМИ.
\end{abstract}

\newpage
\section*{Вступна частина}

\subsection*{Умовні скорочення та визначення}

Терміни та скорочення, що використовуються в документі:

\begin{longtable}{p{4cm}p{10cm}}
\toprule
\textbf{Термін} & \textbf{Значення} \\
\midrule
API & Автоматизовані програмні інтеграції, відомі також як “прикладний інтерфейс програмування – application programming interface” \\
БД & База Даних — Організована колекція даних, яка зберігається та управляється з використанням комп'ютерної системи; забезпечує зберігання, організацію та доступ до даних з метою ефективного використання та обробки \\
БП & Бізнес Процес — Сукупність дій необхідних для здійснення людиною або системою, в результаті яких відбувається досягнення бажаного результату \\
КЕП & Кваліфікований електронний підпис (X.509) \\
РМК & Реєстраційно-моніторингова картка \\
АРМ & Автоматизоване робоче місце користувача \\
ЄСІКС & Єдина судова інформаційно-комунікаційна система \\
BPMN & Business Process Model and Notation (ISO 19510) — стандарт моделювання та оркестрації бізнес-процесів \\
\bottomrule
\end{longtable}

\newpage
\subsection*{Загальні положення}

Технічні вимоги на основі «Концепції ЄСІКС» від 2024 року \cite{esiks}.

\subsection*{Призначення та цілі впровадження}
З метою забезпечення узгодженої, безпечної та ефективної роботи ЄСІКС у сфері документообігу,
а також унеможливлення фрагментації підходів до управління документами,
в межах ЄСІКС є необхідним створення та впровадження централізованого документообігу судочинства:

Реалізація зазначених складових є передумовою побудови цілісної, масштабованої та керованої архітектури ЄСІКС,
зниження витрат на розробку та супровід, забезпечення єдиних стандартів безпеки та користувацького досвіду.

\newpage
\section{Справи}

\subsection{Вимоги до сутностей}

\subsubsection{Адміністративний суд}

\subsubsection{Цивільний суд}

\subsubsection{Господарський суд}

\newpage
\subsection{Вимоги до процесів}

\subsubsection{Первісне подання}

Стадії процесу:

\begin{itemize}
    \item Отримання документів у паперовому вигляді (поштою, особисто або роздруківка)
    \item Отримання документів в електронному вигляді через систему «Електронний суд»
    \item Перевірка цілісності поштового конверту та пакування документів
    \item Перевірка належного оформлення документів (наявність необхідних реквізитів)
    \item Перевірка документів на наявність і відповідність додатків
    \item Перевірка доказів сплати судового збору
    \item Перевірка доказів направлення копій заяви з додатками іншим учасникам справи
    \item Перевірка електронного підпису (КЕП) на документах
    \item Реєстрація документів (створення запису в журналі, проставлення штампу)
    \item Сканування документів та завантаження їх в автоматизовану систему суду (за потреби)
    \item Складання акту про невідповідність документів встановленим вимогам
    \item Створення електронного судового провадження
    \item Долучення документів до електронного провадження
    \item Створення паперової справи (вкладення матеріалів до обгортки)
    \item Виконання авторозподілу справи між суддями
    \item Формування реєстру передачі документів судді для підпису
    \item Друк паперових копій документів, що надійшли в електронному вигляді
    \item Визначення необхідності проведення авторозподілу для конкретного документа
    \item Перевірка правильності реквізитів, особи заявника та його повноважень
    \item Перевірка змістовного характеру документа (відхилення рекламних та недопустимих)
\end{itemize}

\newpage
\subsubsection{Первісний судовий розгляд (підготовчі дії)}

Стадії процесу:

\begin{itemize}
    \item Перевірка відповідності позовної або іншої процесуальної заяви вимогам
    \item Прийняття рішення про відкриття провадження у справі
    \item Залишення позовної або іншої процесуальної заяви без руху
    \item Повернення позовної або іншої процесуальної заяви
    \item Відмова у відкритті провадження у справі
    \item Відмова у поновленні процесуального строку
    \item Залишення позовної або іншої процесуальної заяви без розгляду
    \item Передача справи на розгляд до іншого суду за належністю
    \item Задоволення або відмова у задоволенні заяви про відвід / самовідвід судді
    \item Перевірка на наявність конфлікту інтересів у судді
    \item Поновлення процесуального строку на подання заяви
    \item Розгляд окремих процесуальних заяв (забезпечення позову, витребування доказів тощо)
    \item Відстрочення, розстрочення або звільнення від сплати судового збору
    \item Виклик свідків та витребування доказів
    \item Долучення третіх осіб, які не заявляють самостійних вимог
    \item Визначення виду провадження (загальне або спрощене)
    \item Призначення справи до колегіального розгляду
    \item Постановлення ухвали за результатами підготовчих дій
    \item Внесення ухвали судді до автоматизованої системи суду ДСС/Д3
    \item Направлення судових повідомлень та викликів учасникам справи
    \item Виконання авторозподілу справи (у разі відводу судді)
    \item Прийняття додаткового процесуального рішення без виклику учасників
    \item Направлення запитів до державних органів та установ
\end{itemize}

\newpage
\subsubsection{Первісний судовий розгляд}

Стадії процесу:

\begin{itemize}
    \item Фіксація судового засідання (технічними засобами суду або учасника)
    \item Ведення протоколу судового засідання
    \item Проведення судового засідання в режимі офлайн або онлайн
    \item Закриття підготовчого провадження та призначення справи до судового розгляду по суті
    \item Розгляд справи у підготовчому провадженні або вирішення процесуальних питань
    \item Затвердження мирової угоди або врегулювання спору за участі судді
    \item Зупинення провадження у справі
    \item Залишення позовної заяви без розгляду
    \item Закриття провадження у справі
    \item Передача справи до іншого суду або направлення судового доручення
    \item Задоволення або відмова у задоволенні заяви про відвід / самовідвід судді
    \item Поновлення процесуального строку
    \item Відкладення судового засідання або оголошення перерви
    \item Заміна неналежного відповідача / залучення співвідповідача / третіх осіб
    \item Зміна предмета або підстави позову, збільшення / зменшення позовних вимог
    \item Прийняття додаткового рішення або постановлення окремої ухвали
    \item Встановлення процесуальних строків для вчинення дій учасниками справи
    \item Вирішення інших процесуальних питань відповідно до закону
    \item Постановлення ухвали за результатами розгляду справи
    \item Підготовка документів та внесення їх до автоматизованої системи суду (ДСС/Д3)
    \item Виконання авторозподілу справи (у разі відводу судді)
    \item Направлення судових повідомлень та викликів учасникам справи
\end{itemize}

\newpage
\subsubsection{Розгляд справи по суті}

Стадії процесу:

\begin{itemize}
    \item Відкриття судового засідання по суті та розгляд справи по суті
    \item Ведення протоколу судового засідання
    \item Фіксація судового засідання технічними засобами
    \item Проведення судового засідання в режимі офлайн або онлайн (ВКЗ)
    \item Відкладення судового засідання або оголошення перерви
    \item Повернення справи на підготовче судове засідання
    \item Зупинення провадження у справі
    \item Залишення позовної заяви без розгляду
    \item Закриття провадження у справі
    \item Затвердження мирової угоди
    \item Задоволення або відмова у задоволенні заяви про відвід / самовідвід судді
    \item Вирішення інших процесуальних питань
    \item Винесення судового рішення по суті позову (повністю або частково)
    \item Постановлення ухвали за результатами розгляду
    \item Підготовка та внесення судових документів до автоматизованої системи ДСС/Д3
    \item Формування матеріалів судової справи
    \item Внесення судового рішення до реєстрів та електронних баз даних
    \item Направлення судового рішення учасникам справи та розміщення для публічного доступу
    \item Виконання авторозподілу справи (у разі відводу судді)
\end{itemize}

\newpage
\subsubsection{Виконання судового рішення, судовий контроль}

Стадії процесу:

\begin{itemize}
    \item Ведення протоколу судового засідання
    \item Здійснення фіксації судового засідання
    \item Організація проведення судового засідання в режимі офлайн
    \item Організація проведення судового засідання в режимі онлайн
    \item Організація проведення судового засідання онлайн засобами ВКЗ
    \item Організація проведення судового засідання власними технічними засобами суду
    \item Видача виконавчого документа (у паперовій або електронній формі через ЄСІТС)
    \item Внесення відомостей до ЄДР виконавчих документів ЄДРВД
    \item Внесення відомостей до ЄДР РНОККП, ФОП, ЄДРРП, тощо
    \item Здійснення судового контролю за виконанням судових рішень
    \item Вирішення скарг на рішення, дії або бездіяльність органів виконання судових рішень
    \item Розгляд звіту відповідача про виконання судового рішення
    \item Здійснення перевірки виконання вказівок, що містяться в окремій ухвалі суду
    \item Вирішення питання про поворот виконання судового рішення
    \item Виправлення помилок у виконавчому документі
    \item Відстрочення або розстрочення виконання судового рішення
    \item Зупинення виконання судового рішення
    \item Визнання виконавчого документа таким, що не підлягає виконанню
    \item Видача дубліката виконавчого документа
    \item Вирішення інших процесуальних питань, пов’язаних із виконанням судових рішень
    \item Постановлення ухвали суду за результатами розгляду процесуальних питань
    \item Підготовка документів та внесення їх до ДСС/Д3
    \item Перевірка на наявність конфлікту інтересів судді
    \item Розгляд заяви про відвід / самовідвід судді
    \item Авторозподіл справи
    \item Направлення справи до іншого суду для вирішення питання про відвід
\end{itemize}

\newpage
\subsubsection{Подання апеляційної, касаційної скарги}

Стадії процесу:

\begin{itemize}
    \item Отримання апеляційної або касаційної скарги в паперовому вигляді
    \item Отримання апеляційної або касаційної скарги в електронному вигляді через «Електронний суд»
    \item Перевірка цілісності поштового конверту та пакування документів
    \item Перевірка належного оформлення скарги (наявність необхідних реквізитів)
    \item Перевірка документів на наявність і відповідність додатків
    \item Перевірка доказів сплати судового збору
    \item Перевірка доказів направлення копій скарги з додатками іншим учасникам справи
    \item Перевірка електронного підпису (КЕП) на документах
    \item Складання акту про невідповідність документів встановленим вимогам
    \item Реєстрація скарги (створення запису в журналі, проставлення штампу)
    \item Сканування документів та завантаження їх в автоматизовану систему суду (за потреби)
    \item Створення електронного провадження
    \item Долучення документів до електронного провадження
    \item Створення паперової справи (вкладення матеріалів до обгортки)
    \item Формування реєстру передачі документів судді для підпису
    \item Виконання авторозподілу скарги між суддями
    \item Підготовка та внесення документів до системи документообігу суду (ДСС/Д3)
    \item Друк паперових копій документів, що надійшли в електронному вигляді
\end{itemize}

\newpage
\subsubsection{Апеляційний розгляд}

Стадії процесу:

\begin{itemize}
    \item Отримання апеляційної скарги або матеріалів справи з суду першої інстанції
    \item Перевірка апеляційної скарги на відповідність вимогам закону
    \item Вирішення питання про відкриття апеляційного провадження
    \item Залишення апеляційної скарги без руху
    \item Повернення апеляційної скарги
    \item Відмова у відкритті апеляційного провадження
    \item Відкриття апеляційного провадження
    \item Забезпечення повідомлення (виклику) учасників процесу
    \item Підготовка справи до апеляційного розгляду
    \item Проведення судового засідання в режимі офлайн або онлайн
    \item Ведення протоколу судового засідання
    \item Фіксація судового засідання технічними засобами
    \item Розгляд заяви про відвід (самовідвід) судді
    \item Відкладення розгляду справи або оголошення перерви
    \item Вирішення інших процесуальних питань
    \item Розгляд апеляційної скарги по суті
    \item Залишення рішення суду першої інстанції без змін
    \item Скасування рішення суду першої інстанції повністю або частково
    \item Зміна рішення суду першої інстанції
    \item Скасування рішення та направлення справи на новий розгляд
    \item Закриття апеляційного провадження
    \item Постановлення ухвали або постанови апеляційного суду
    \item Підготовка та внесення судових документів до ДСС/Д3
    \item Внесення судового рішення до реєстрів та електронних баз даних
    \item Надсилання судового рішення учасникам справи та забезпечення публічного доступу
    \item Повернення матеріалів справи до суду першої інстанції
\end{itemize}

\newpage
\subsubsection{Касаційний розгляд}

Стадії процесу:

\begin{itemize}
    \item Отримання касаційної скарги або матеріалів справи з нижчих інстанцій
    \item Перевірка касаційної скарги на відповідність вимогам закону
    \item Вирішення питання про відкриття касаційного провадження
    \item Залишення касаційної скарги без руху
    \item Повернення касаційної скарги
    \item Відмова у відкритті касаційного провадження
    \item Відкриття касаційного провадження
    \item Забезпечення повідомлення (виклику) учасників процесу
    \item Підготовка справи до касаційного розгляду
    \item Проведення судового засідання в режимі офлайн або онлайн
    \item Ведення протоколу судового засідання
    \item Фіксація судового засідання технічними засобами
    \item Розгляд заяви про відвід (самовідвід) судді
    \item Відкладення розгляду справи або оголошення перерви
    \item Вирішення інших процесуальних питань
    \item Розгляд касаційної скарги по суті
    \item Залишення судових рішень без змін
    \item Скасування судових рішень повністю або частково
    \item Зміна судових рішень
    \item Скасування рішень та направлення справи на новий розгляд
    \item Закриття касаційного провадження
    \item Постановлення ухвали або постанови касаційного суду
    \item Підготовка та внесення судових документів до ДСС/Д3
    \item Внесення судового рішення до реєстрів та електронних баз даних
    \item Надсилання судового рішення учасникам справи та забезпечення публічного доступу
    \item Повернення матеріалів справи до суду нижчої інстанції
\end{itemize}

\newpage
\subsubsection{Розгляд справи Верховним судом}

Стадії процесу:

\begin{itemize}
    \item Отримання матеріалів справи для розгляду Верховним Судом
    \item Перевірка підстав для розгляду справи палатами або Великою Палатою
    \item Вирішення питання про передачу справи на розгляд палатами або Великою Палатою
    \item Постановлення ухвали про передачу справи на розгляд палати або Великої Палати
    \item Розгляд справи колегією суддів Верховного Суду
    \item Розгляд справи палатою Верховного Суду
    \item Розгляд справи об’єднаною палатою Верховного Суду
    \item Розгляд справи Великою Палатою Верховного Суду
    \item Повернення справи колегії, палаті або об’єднаній палаті для продовження розгляду
    \item Постановлення ухвали за результатами розгляду
    \item Підготовка та внесення судових документів до ДСС/Д3
    \item Внесення судового рішення до реєстрів та електронних баз даних
    \item Надсилання судового рішення учасникам справи та забезпечення публічного доступу
\end{itemize}

\newpage
\section{Провадження}

\subsection{Вимоги до сутностей}

\subsubsection{Слідчий}

\subsubsection{Судовий слідчий}

\subsubsection{Прокурор}

\subsubsection{Адвокат}

\newpage
\subsection{Вимоги до процесів}

\subsubsection{Судовий розгляд під час досудового розслідування}

Стадії процесу:

\begin{itemize}
    \item Повідомити учасників кримінального провадження та інших осіб про судове засідання
    \item Повідомити слідчого, дізнавача або прокурора про судове засідання
    \item Підготувати документи та внести їх до електронних систем суду (ДСС/Д3)
    \item Фіксувати судове засідання
    \item Вести протокол судового засідання
    \item Провести судове засідання в режимі онлайн (ВКЗ) або офлайн
    \item Розглянути клопотання про тимчасовий доступ до речей і документів
    \item Розглянути клопотання про арешт майна або скасування арешту
    \item Розглянути клопотання про застосування запобіжного заходу
    \item Розглянути клопотання про дозвіл на затримання з метою приводу
    \item Розглянути клопотання про привід
    \item Розглянути клопотання про накладення грошового стягнення
    \item Розглянути клопотання про відсторонення від посади
    \item Розглянути клопотання про огляд з метою безпосереднього дослідження доказів
    \item Перевірити наявність конфлікту інтересів
    \item Розглянути заяву про відвід або самовідвід судді
    \item Вирішити питання самовідводу судді
    \item Передати справу на авторозподіл для нового складу суду (у разі відводу)
    \item Направити справу до іншого суду (якщо неможливо сформувати новий склад)
    \item Постановити ухвалу за результатами розгляду клопотання
    \item Оформити справу та передати її за підсудністю
    \item Роз’яснити судове рішення за клопотанням учасника або органу виконання
    \item Виправити описки чи технічні помилки в судовому рішенні
\end{itemize}

\newpage
\subsubsection{Скарги в рамках досудового розслідування}

Стадії процесу:

\begin{itemize}
    \item Повідомити особу, яка подала скаргу, та третіх осіб про судове засідання
    \item Повідомити слідчого, дізнавача або прокурора про судове засідання
    \item Підготувати документи та внести їх до електронних систем суду (ДСС/Д3)
    \item Фіксувати судове засідання
    \item Вести протокол судового засідання
    \item Провести судове засідання в режимі онлайн (ВКЗ) або офлайн
    \item Розглянути скаргу на бездіяльність щодо невнесення відомостей до ЄРДР
    \item Розглянути скаргу на рішення про зупинення досудового розслідування
    \item Розглянути скаргу на рішення про закриття кримінального провадження
    \item Розглянути скаргу на рішення прокурора про закриття провадження
    \item Розглянути інші скарги, передбачені главою 26 КПК України
    \item Здійснити розгляд скарги за участю сторін
    \item Перевірити наявність конфлікту інтересів
    \item Розглянути заяву про відвід або самовідвід слідчого судді
    \item Вирішити питання самовідводу судді
    \item Передати справу на авторозподіл для нового складу суду
    \item Направити справу до іншого суду (якщо неможливо сформувати склад)
    \item Повернути скаргу заявнику
    \item Відмовити у задоволенні скарги повністю або частково
    \item Задовольнити скаргу повністю або частково
    \item Відмовити у відкритті провадження за скаргою
    \item Постановити ухвалу за результатами розгляду скарги
    \item Оформити справу та передати її за підсудністю
\end{itemize}

\newpage
\subsubsection{Клопотання в рамках досудового розслідування}

Стадії процесу:

\begin{itemize}
    \item Повідомити особу, яка подала клопотання, та третіх осіб про судове засідання
    \item Повідомити слідчого, дізнавача або прокурора про судове засідання
    \item Підготувати документи та внести їх до електронних систем суду (ДСС/Д3)
    \item Фіксувати судове засідання
    \item Вести протокол судового засідання
    \item Провести судове засідання в режимі онлайн (ВКЗ) або офлайн
    \item Розглянути клопотання про накладення грошового стягнення
    \item Розглянути клопотання про обмеження у користуванні спеціальним правом
    \item Розглянути клопотання про відсторонення від посади
    \item Розглянути клопотання про тимчасовий доступ до речей і документів
    \item Розглянути клопотання про накладення арешту на майно
    \item Розглянути клопотання про скасування арешту майна
    \item Розглянути клопотання про привід
    \item Розглянути інші клопотання, передбачені КПК України
    \item Перевірити наявність конфлікту інтересів
    \item Розглянути заяву про відвід або самовідвід слідчого судді
    \item Вирішити питання самовідводу судді
    \item Передати справу на авторозподіл для нового складу суду
    \item Направити справу до іншого суду (якщо неможливо сформувати склад)
    \item Повернути клопотання заявнику
    \item Відмовити у задоволенні клопотання повністю або частково
    \item Задовольнити клопотання повністю або частково
    \item Залишити клопотання без розгляду
    \item Відмовити у відкритті провадження за клопотанням
    \item Постановити ухвалу за результатами розгляду клопотання
    \item Оформити справу та передати її за підсудністю
\end{itemize}

\newpage
\subsubsection{Судовий розгляд кримінальних проступків}

Стадії процесу:

\begin{itemize}
    \item Повідомити учасників кримінального провадження та інших осіб про судове засідання
    \item Повідомити слідчого, дізнавача або прокурора про судове засідання
    \item Підготувати документи та внести їх до електронних систем суду (ДСС/Д3)
    \item Фіксувати судове засідання
    \item Вести протокол судового засідання
    \item Провести судове засідання в режимі онлайн (ВКЗ) або офлайн
    \item Викликати учасників кримінального провадження на судовий розгляд
    \item Провести судовий розгляд кримінального проступку в судовому засіданні
    \item Провести судовий розгляд без проведення судового засідання (у передбачених випадках)
    \item Призначити судовий розгляд по суті
    \item Ухвалити вирок (обвинувальний або виправдувальний)
    \item Перевірити наявність конфлікту інтересів
    \item Розглянути заяву про відвід або самовідвід судді
    \item Вирішити питання самовідводу судді
    \item Передати справу на авторозподіл для нового складу суду
    \item Направити справу до іншого суду (якщо неможливо сформувати склад)
    \item Постановити ухвалу за результатами розгляду
    \item Оформити справу та передати її за підсудністю
\end{itemize}

\newpage
\subsubsection{Підготовче судове засідання}

Стадії процесу:

\begin{itemize}
    \item Повідомити учасників кримінального провадження та інших осіб про підготовче засідання
    \item Повідомити слідчого, дізнавача або прокурора про підготовче засідання
    \item Підготувати документи та внести їх до електронних систем суду (ДСС/Д3)
    \item Здійснити авторозподіл справи
    \item Фіксувати підготовче судове засідання
    \item Вести протокол підготовчого судового засідання
    \item Провести підготовче судове засідання в режимі онлайн (ВКЗ) або офлайн
    \item Розглянути клопотання про затвердження угоди про визнання винуватості
    \item Розглянути клопотання про затвердження угоди про примирення
    \item Розглянути клопотання про застосування примусових заходів медичного чи виховного характеру
    \item Розглянути клопотання про звільнення особи від кримінальної відповідальності
    \item Розглянути клопотання про примирення в порядку приватного обвинувачення
    \item Розглянути обвинувальний акт та клопотання прокурора про спрощене провадження
    \item Розглянути клопотання про закриття кримінального провадження
    \item Вирішити питання щодо міри запобіжного заходу або його продовження
    \item Вирішити питання про складання досудової доповіді
    \item Розглянути клопотання учасників про витребування доказів, призначення експертизи тощо
    \item Вирішити питання про відвід або самовідвід судді, прокурора, захисника чи інших учасників
    \item Вирішити питання про видалення порушника із зали судового засідання
    \item Замінити присяжних запасними (за потреби)
    \item Провести процесуальні дії в режимі відеоконференції
    \item Зупинити судове провадження
    \item Об’єднати або виділити матеріали кримінального провадження
    \item Повернути обвинувальний акт або клопотання прокурору через невідповідність вимогам
    \item Направити справу до іншого суду для визначення підсудності
    \item Призначити судовий розгляд по суті
    \item Постановити ухвалу за результатами підготовчого засідання
    \item Роз’яснити судове рішення за клопотанням учасника або органу виконання
    \item Розглянути заяву про відвід/самовідвід поточним складом суду
    \item Передати справу на авторозподіл для нового складу суду (у разі відводу)
\end{itemize}

\newpage
\subsubsection{Судовий розгляд}

Стадії процесу:

\begin{itemize}
    \item Отримати процесуальні документи в електронному вигляді через ЄСІТС («Електронний суд»)
    \item Отримати документи в паперовому вигляді (поштою, нарочним або роздруківка)
    \item Перевірити цілісність поштового конверту та документів
    \item Перевірити документи на наявність та відповідність додатків
    \item Перевірити докази направлення цивільного позову або процесуального документа учасникам
    \item Перевірити правильність реквізитів, належність оформлення та особу заявника
    \item Перевірити зміст документа (відхилити рекламні та інші невідповідні документи)
    \item Перевірити КЕП (кваліфікований електронний підпис)
    \item Скласти акт про невідповідність документів умовам прийняття
    \item Зареєструвати документ у відповідному журналі та проставити реєстраційну позначку
    \item Створити електронне провадження (або справу в паперовому вигляді)
    \item Долучити документ до електронного провадження
    \item Сформувати реєстр передачі документів судді
    \item Передати документи судді після реєстрації та перевірки
    \item Здійснити авторозподіл (у разі потреби)
    \item Сканувати паперові документи та завантажити їх в систему (якщо надійшли не через ЄСІТС)
\end{itemize}


\newpage
\subsubsection{Виконання судового рішення, судовий контроль}

Стадії процесу:

\begin{itemize}
    \item Направити судове рішення до виконання (видача виконавчого документа)
    \item Надіслати розпорядження та копію рішення органу або установі виконання
    \item Розглянути клопотання (подання) про вирішення питань, пов’язаних з виконанням вироку
    \item Розглянути питання про умовно-дострокове звільнення від відбування покарання
    \item Розглянути питання про заміну невідбутої частини покарання більш м’яким
    \item Розглянути питання про звільнення від відбування покарання жінок з дітьми до 3 р. і вагітних
    \item Розглянути питання про відстрочку виконання вироку
    \item Розглянути питання про зняття судимості
    \item Розглянути інші питання, пов’язані з виконанням судового рішення (ст. 537 КПК)
    \item Здійснити судовий контроль за виконанням судового рішення
    \item Перевірити наявність конфлікту інтересів
    \item Розглянути заяву про відвід або самовідвід судді
    \item Вирішити питання самовідводу судді
    \item Передати справу на авторозподіл для нового складу суду
    \item Направити справу до іншого суду (якщо неможливо сформувати склад)
    \item Постановити ухвалу за результатами розгляду клопотання/подання
    \item Підготувати документи та внести їх до електронних систем суду (ДСС/Д3)
    \item Оформити справу та передати її за підсудністю
\end{itemize}


\newpage
\subsubsection{Подання апеляційної, касаційної скарги}

Стадії процесу:

\begin{itemize}
    \item Отримати апеляційну або касаційну скаргу в електронному вигляді через ЄСІТС
    \item Отримати скаргу в паперовому вигляді (поштою, нарочним або роздруківка)
    \item Перевірити цілісність поштового конверту та документів
    \item Перевірити документи на наявність та відповідність додатків
    \item Перевірити правильність реквізитів, належність оформлення та повноваження заявника
    \item Перевірити зміст документа (відхилити рекламні та інші невідповідні документи)
    \item Перевірити КЕП (кваліфікований електронний підпис)
    \item Скласти акт про невідповідність документів умовам прийняття
    \item Зареєструвати скаргу у відповідному журналі та проставити реєстраційну позначку
    \item Створити електронне провадження (або справу в паперовому вигляді)
    \item Долучити скаргу до електронного провадження
    \item Сформувати реєстр передачі документів судді
    \item Передати скаргу судді після реєстрації та перевірки
    \item Здійснити авторозподіл (у разі потреби)
    \item Сканувати паперові документи та завантажити їх в систему (якщо надійшли не через ЄСІТС)
\end{itemize}

\newpage
\subsubsection{Апеляційний розгляд}

Стадії процесу:

\begin{itemize}
    \item Отримати матеріали справи з суду першої інстанції
    \item Перевірити повноту матеріалів апеляційної скарги
    \item Здійснити підготовку до апеляційного розгляду
    \item Забезпечити повідомлення та виклик учасників кримінального провадження
    \item Здійснити авторозподіл справи (у разі потреби)
    \item Перевірити наявність конфлікту інтересів
    \item Розглянути заяву про відвід або самовідвід судді
    \item Вирішити питання самовідводу судді
    \item Передати справу на авторозподіл для нового складу суду
    \item Провести апеляційний розгляд у судовому засіданні
    \item Провести апеляційний розгляд у порядку письмового провадження
    \item Постановити ухвалу за результатами апеляційного розгляду
    \item Залишити вирок (ухвалу) без змін
    \item Змінити вирок (ухвалу)
    \item Скасувати вирок (ухвалу) повністю або частково та ухвалити нове рішення
    \item Скасувати вирок (ухвалу) і закрити кримінальне провадження
    \item Скасувати вирок (ухвалу) і призначити новий розгляд у суді першої інстанції
    \item Підготувати документи та внести їх до електронних систем суду (ДСС/Д3)
    \item Оформити справу та повернути її до суду першої інстанції
    \item Внести судове рішення до відповідних реєстрів
\end{itemize}

\newpage
\subsubsection{Касаційний розгляд}

Стадії процесу:

\begin{itemize}
    \item Отримати матеріали касаційної скарги та справи з нижчих інстанцій
    \item Здійснити підготовку до касаційного розгляду
    \item Забезпечити повідомлення та виклик учасників кримінального провадження
    \item Перевірити наявність конфлікту інтересів
    \item Розглянути заяву про відвід або самовідвід судді
    \item Вирішити питання самовідводу судді
    \item Передати справу на авторозподіл для нового складу суду
    \item Вирішити питання про відкриття касаційного провадження
    \item Відкрити касаційне провадження
    \item Відмовити у відкритті касаційного провадження
    \item Залишити касаційну скаргу без руху
    \item Повернути касаційну скаргу
    \item Здійснити касаційний розгляд (у судовому засіданні або в порядку письмового провадження)
    \item Постановити ухвалу за результатами касаційного розгляду
    \item Залишити судове рішення без змін
    \item Змінити судове рішення
    \item Скасувати судове рішення повністю або частково та ухвалити нове рішення
    \item Скасувати судове рішення і закрити кримінальне провадження
    \item Скасувати судове рішення і призначити новий розгляд у суді першої або апеляційної інстанції
    \item Підготувати документи та внести їх до електронних систем суду (ДСС/Д3)
    \item Оформити справу та повернути її до суду першої або апеляційної інстанції
    \item Внести судове рішення до відповідних реєстрів
\end{itemize}

\newpage
\subsubsection{Розгляд справи Верховним судом}

Стадії процесу:

\begin{itemize}
    \item Отримати справу для розгляду палатою, об’єднаною палатою або Великою Палатою Верховного Суду
    \item Вирішити питання про наявність підстав для розгляду справи палатою, об’єднаною палатою або Великою Палатою
    \item Передати справу на розгляд палати Верховного Суду
    \item Передати справу на розгляд об’єднаної палати Верховного Суду
    \item Передати справу на розгляд Великої Палати Верховного Суду
    \item Розглянути справу колегією суддів (палатою / об’єднаною палатою / Великою Палатою)
    \item Постановити ухвалу за результатами розгляду
    \item Повернути справу до попередньої інстанції або передати за підсудністю
    \item Підготувати документи та внести їх до електронних систем суду (ДСС/Д3)
    \item Оформити справу та повернути її до суду першої або апеляційної інстанції
\end{itemize}

\addtocontents{toc}{\protect\newpage}
\newpage
\section{Судові засідання і їх рішення}

\newpage
\section{Авторозподіл справ}

\subsection{Вимоги до процесів}

\subsubsection{Авторозподілення справи між суддями}

\newpage
\section{Висновок}

\begin{thebibliography}{9}

\bibitem{ecoz:rehab} Міністерство охорони здоров’я України. \newblock Технічні вимоги Реабілітація 2.0. \newblock Версія 08.08.2024.
\bibitem{dstu:3973} ДСТУ 3973-2000. \newblock Система розроблення та постачання продукції на виробництво. Правила виконання науково-дослідних робіт. Загальні положення.
\bibitem{dstu:3974} ДСТУ 3974-2000. \newblock Система розроблення та постачання продукції на виробництво. Правила виконання науково-конструкторських робіт. Загальні положення.
\bibitem{dstu:42010} ДСТУ 42010:2018. \newblock Інженерія систем і програмних засобів. Опис архітектури.
\bibitem{law:info-protection} Закон України «Про захист інформації в інформаційно-телекомунікаційних системах».
\bibitem{law:personal-data} Закон України «Про захист персональних даних».
\bibitem{law:cybersecurity} Закон України «Про основні засади забезпечення кібербезпеки України».
\bibitem{owasp:top10} OWASP Foundation. \newblock OWASP Top 10: The Ten Most Critical Web Application Security Risks. \newblock \url{https://owasp.org/www-project-top-ten/}.
\bibitem{nist:sp800-57} NIST Special Publication 800-57 Part 1 Revision 5. \newblock Recommendation for Key Management.
\bibitem{iso:42010} ISO/IEC/IEEE 42010:2022. \newblock Systems and software engineering — Architecture description.
\bibitem{zachman} Zachman, J.A. \newblock A Framework for Information Systems Architecture. \newblock IBM Systems Journal, 1987.

\end{thebibliography}

\end{document}
